% =========================================================
% capitulos/06_Discusion.tex
% Placeholder con las herramientas que se van a usar en la
% tesis: tablas con el estilo azul, pseudocódigo y ecuaciones
% numeradas / sin numerar de distintos tipos.
% =========================================================

\chapter{Resultados}

\label{chap:Resultados}


\section{Resultados del problema de muestras pequeñas (SSS)}

\subsection{Caso 1}

\begin{figure}[htbp]
	\centering
	
	% Primera fila
	\begin{subfigure}[b]{0.3\textwidth}
		\centering
		\includegraphics[width=\textwidth]{caso_1_1_lda_clasico.png}
	\end{subfigure}
	\hfill
	\begin{subfigure}[b]{0.3\textwidth}
		\centering
		\includegraphics[width=\textwidth]{caso_1_2_lda_mass.png}
	\end{subfigure}
	\hfill
	\begin{subfigure}[b]{0.3\textwidth}
		\centering
		\includegraphics[width=\textwidth]{caso_1_3_kfda.png}
	\end{subfigure}
		
	\caption{Comparativa de los espacios latentes generados por LDA, LDA MASS y KFDA ante
		el impacto de la alta dimensionalidad (caso 1).}
	\label{fig:R1}
\end{figure}


La figura \ref{fig:R1} muestra las proyecciones obtenidas por los tres métodos en un escenario ideal
($n>p$). Se observa que la estructura geométrica del mapeo lograda por el LDA clásico y
LDA MASS es sumamente similar, diferenciándose principalmente por la orientación de los
datos y la escala de los ejes. Por su parte, la proyección del KFDA presenta un
comportamiento geométrico distinto: las clases quedan visualmente más alejadas entre sí y
altamente compactas, a diferencia de los modelos LDA que exhiben una mayor dispersión
intra-clase.

\subsection{Caso 2}

\begin{figure}[htbp]
	\centering
	
	% Primera fila
	\begin{subfigure}[b]{0.3\textwidth}
		\centering
		\includegraphics[width=\textwidth]{caso_2_1_lda_clasico.png}
	\end{subfigure}
	\hfill
	\begin{subfigure}[b]{0.3\textwidth}
		\centering
		\includegraphics[width=\textwidth]{caso_2_2_lda_mass.png}
	\end{subfigure}
	\hfill
	\begin{subfigure}[b]{0.3\textwidth}
		\centering
		\includegraphics[width=\textwidth]{caso_2_3_kfda.png}
	\end{subfigure}
	
	\caption{Comparativa de los espacios latentes generados por LDA, LDA MASS y KFDA ante
		el impacto de la alta dimensionalidad (caso 2).}
	\label{fig:R2}
\end{figure}

La Figura \ref{fig:R2} muestra las proyecciones obtenidas por los tres métodos en un escenario
de mayor complejidad donde ($n<p$ ). Se observa que pese a ser más variables que
observaciones, la estructura geométrica del mapeo lograda por el LDA clásico y LDA MASS
es similar, diferenciándose principalmente por la orientación de los datos y la escala de los
ejes nuevamente. Por su parte, la proyección del KFDA genera un espacio ultra compacto
donde la dispersión intra-clase es mínima y la dispersión entre-clases es máxima en un
espacio muy reducido.

\subsection{Caso 3}

\begin{figure}[htbp]
	\centering
	
	% Primera fila
	\begin{subfigure}[b]{0.3\textwidth}
		\centering
		\includegraphics[width=\textwidth]{caso_3_1_lda_clasico.png}
	\end{subfigure}
	\hfill
	\begin{subfigure}[b]{0.3\textwidth}
		\centering
		\includegraphics[width=\textwidth]{caso_3_2_lda_mass.png}
	\end{subfigure}
	\hfill
	\begin{subfigure}[b]{0.3\textwidth}
		\centering
		\includegraphics[width=\textwidth]{caso_3_3_kfda.png}
	\end{subfigure}
	
	\caption{Comparativa de los espacios latentes generados por LDA, LDA MASS y KFDA ante
		el impacto de la alta dimensionalidad (caso 3).}
	\label{fig:R3}
\end{figure}

La Figura \ref{fig:R3} muestra las proyecciones obtenidas por los tres métodos en un escenario
de mayor complejidad donde ($n<p$). Pese a que la cantidad de observaciones por grupo es
menor a la cantidad de variables presentes, los modelos aún son capaces de ofrecer
proyecciones de los datos. La estructura geométrica del mapeo lograda por el LDA clásico y
LDA MASS es similar, diferenciándose principalmente por la orientación de los datos y la
escala de los ejes nuevamente. Por su parte, la proyección del KFDA vuelve a generar una
proyección donde la dispersión intra-clase es mínima y la dispersión entre clase es máxima,
diferenciándose de los comportamientos mostrados por los demás métodos.

\subsection{Caso 4}

\begin{figure}[htbp]
	\centering
	
	% Primera fila
	\begin{subfigure}[b]{0.3\textwidth}
		\centering
		\includegraphics[width=\textwidth]{caso_4_1_lda_clasico.png}
	\end{subfigure}
	\hfill
	\begin{subfigure}[b]{0.3\textwidth}
		\centering
		\includegraphics[width=\textwidth]{caso_4_2_lda_mass.png}
	\end{subfigure}
	\hfill
	\begin{subfigure}[b]{0.3\textwidth}
		\centering
		\includegraphics[width=\textwidth]{caso_4_3_kfda.png}
	\end{subfigure}
	
	\caption{Comparativa de los espacios latentes generados por LDA, LDA MASS y KFDA ante
		el impacto de la alta dimensionalidad (caso 4).}
	\label{fig:R4}
\end{figure}

La Figura \ref{fig:R4} ilustra el impacto de una condición extrema del problema de tamaño de
muestra pequeño ($p>n$). Ante esta desproporción severa, la matriz de covarianza se vuelve
singular, provocando que el LDA clásico resulte matemáticamente incomputable. No
obstante, las variantes LDA MASS y KFDA logran converger y generar proyecciones del
espacio latente. La proyección obtenida por LDA MASS permite distinguir la separación de
los grupos con una elevada dispersión intra-clase. Por su parte, el mapeo de KFDA induce
una compresión máxima de la varianza intra-clase, colapsando visualmente las observaciones
en coordenadas singulares; si bien esto exhibe una separación inter-clase total, el
comportamiento sugiere fuertemente un sobreajuste geométrico en el espacio transformado.

\subsection{Caso 5}

\begin{figure}[htbp]
	\centering
	
	% Primera fila
	\begin{subfigure}[b]{0.3\textwidth}
		\centering
		\includegraphics[width=\textwidth]{caso_5_1_lda_clasico.png}
	\end{subfigure}
	\hfill
	\begin{subfigure}[b]{0.3\textwidth}
		\centering
		\includegraphics[width=\textwidth]{caso_5_2_lda_mass.png}
	\end{subfigure}
	\hfill
	\begin{subfigure}[b]{0.3\textwidth}
		\centering
		\includegraphics[width=\textwidth]{caso_5_3_kfda.png}
	\end{subfigure}
	
	\caption{Comparativa de los espacios latentes generados por LDA, LDA MASS y KFDA ante
		el impacto de la alta dimensionalidad (caso 5).}
	\label{fig:R5}
\end{figure}

La Figura \ref{fig:R5} ilustra el impacto de una condición extrema del problema de tamaño de
muestra pequeño ($p>n$) más agresivo que el del caso 4. La matriz de covarianza se vuelve
singular, provocando que el LDA clásico resulte matemáticamente incomputable. No
obstante, las variantes LDA MASS y KFDA logran converger y generar proyecciones del
espacio latente. La proyección obtenida por LDA MASS permite distinguir la separación de
los grupos, pero con un cierto grado de traslape entre algunas observaciones de ellos, además
de tener una elevada dispersión intra-clase. Por su parte, el mapeo de KFDA induce una
compresión máxima de la varianza intra-clase y maximizando la separación inter-clase. Sin
embargo, esto puede deberse a un severo problema de sobre ajuste.

%%%%%%%%%%   RECONSTRUIR POR CASOS %%%%%%%%%%%%%%%%%%%%%%%%%%%%%%%%%

\newpage

\section{Resultados por conjunto de datos}

En esta sección se encuentran desarrollados los resultados correspondientes a la exploración de cada unos de los conjuntos de datos, sus proyecciones de baja dimensionalidad obtenidas mediante LDA, KFDA y GNLPDA ademas de los resultados comparativos de rendimientos en cuanto a clasificación. 

\subsection{Palmer Penguins}

El conjunto de datos pingüinos Palmer fue originalmente introducido por K. B. Gorman, T. D. Williams y W. R. Fraser en el 2014 \cite{gorman2014ecological}. El conjunto consiste en 342 observaciones distribuidas desigualmente entre tres especies de pingüinos: Adelie ($44.2\%$), Chinstrap($19.8\%$) y Gentoo ($36\%$). Cada observación consta de cuatro variables continuas numéricas:   largo del pico en milímetros, ancho del pico en milímetros, largo de la aleta en milímetros y la masa corporal en gramos.

Desde la perspectiva del reconocimiento de patrones, el conjunto Palmer Penguins presenta propiedades que lo hacen particularmente útil para la evaluación experimental de métodos de clasificación y reducción de dimensionalidad. Los análisis exploratorios reportados en la literatura indican que las observaciones correspondientes a la especie Gentoo tienden a formar un grupo relativamente diferenciado, mientras que las especies Adelie y Chinstrap exhiben cierto grado de solapamiento en el espacio de características, dando lugar a clases con distintos niveles de separabilidad \cite{horst2022palmer}. Adicionalmente, la distribución desigual de observaciones entre especies introduce un escenario de desbalance de clases que permite examinar la sensibilidad y robustez de los métodos analizados frente a distribuciones no uniformes. En conjunto, estas características hacen de Palmer Penguins un caso de estudio pertinente para analizar la capacidad discriminativa de técnicas de proyección y clasificación, así como su comportamiento ante conjuntos de datos que combinan separabilidad parcial y desbalance entre clases.


A continuación se muestra mediante gráficos de dispersión el comportamiento de las observaciones por especies en pares de variables:

\begin{figure}[htbp]
	\centering
	\includegraphics[width=0.5\textwidth]{1DES.png}
	\caption{Gráficos de dispersión por pares de Palmer Penguins.}
	\label{fig:D1}
\end{figure}

En el gráfico \ref{fig:D1} se observa el comportamiento por pares de variables para el conjunto de pinguinos Palmer. Se observa principalmente que lkos grupos se alcanzan a diferenciar bien entre la longitud del pico y el ancho del pico pero mantienen un grado de traslape principalmente observable entre dos especies al observar la dispersión que se forma entre el ancho del pico y la masa corporal.

A continuación se obtuvieron las proyecciones visuales mediante la reducción de características con LDA, KFDA y GNLPDA.

\begin{figure}[htbp]
	\centering
	
	% Primera fila
	\begin{subfigure}[b]{0.3\textwidth}
		\centering
		\includegraphics[width=\textwidth]{mejor_escenario_1.2.penguins_1_mass_2d.png}
	\end{subfigure}
	\hfill
	\begin{subfigure}[b]{0.3\textwidth}
		\centering
		\includegraphics[width=\textwidth]{mejor_escenario_1.2.penguins_2_kfda_2d.png}
	\end{subfigure}
	\hfill
	\begin{subfigure}[b]{0.3\textwidth}
		\centering
		\includegraphics[width=\textwidth]{mejor_escenario_1.2.penguins_3_gnlpda_2d.png}
	\end{subfigure}
	
	\caption{Comparativa geométrica de los espacios de proyección para el conjunto de datos
		Palmer Penguins.}
	\label{fig:R6}
\end{figure}

La Figura \ref{fig:R6} presenta los espacios latentes generados por los tres enfoques evaluados
con capacidad de proyectar. Dada la baja complejidad geométrica del conjunto Palmer
Penguines, se advierte que la aproximación lineal (LDA MASS) logra una separación de
clases adecuada. Simultáneamente, el mapeo no lineal implícito (KFDA) maximiza de forma
esperada la distancia inter-clase, consolidando agrupaciones visualmente compactas. El
hallazgo más destacable reside en la proyección obtenida por el método propuesto: GNLPDA no solo iguala la capacidad de separación óptima del kernel gaussiano y del LDA MASS, sino que alcanza este rendimiento filtrando y seleccionando explícitamente solo 3 variables
del espacio de características expandido, preservando así la interpretabilidad y trazabilidad
del modelo obtenido.


\begin{figure}[htbp]
	\centering
	\includegraphics[width=0.8\textwidth]{Boxplot_1.2.penguins_Accuracy.png}
	\caption{Distribución empírica del rendimiento predictivo para el conjunto de datos
		Penguins.}
	\label{fig:R19}
\end{figure}

La Figura \ref{fig:R19} presenta la distribución del rendimiento a lo largo de las ejecuciones
experimentales. Los resultados revelan un escenario de alta separabilidad donde la mayoría
de los algoritmos incluyendo enfoques lineales (LDA, QDA) y máquinas de vectores de
soporte (SVM) convergen en medianas superiores a 0.97 con rangos intercuartílicos
estrechos. La excepción estadística es el modelo KFDA, el cual exhibe una caída significativa
en su rendimiento central y una varianza pronunciada con múltiples valores atípicos
inferiores a 0.95. Por su parte, el algoritmo GNLPDA demuestra un comportamiento 
empírico robusto y competitivo; consolida una mediana equiparable a los mejores enfoques
evaluados y alcanza la exactitud perfecta (1.0) en el límite superior de su distribución. 

\begin{figure}[htbp]
	\centering
	\includegraphics[width=0.8\textwidth]{CDplot_Nemenyi_1.2.penguins_long_graficos_todos.png}
	\caption{Análisis de significancia estadística post-hoc mediante distancias críticas de
		Nemenyi para el conjunto de datos Palmer penguins.}
	\label{fig:R20}
\end{figure}


Con el objetivo de validar formalmente las diferencias observadas en el rendimiento
predictivo, se aplicó la prueba de rangos de Friedman, la cual arrojó un estadístico $\chi^2 = 48.175,p<0.001$, proporcionando evidencia suficiente para rechazar la hipótesis de equivalencia entre los clasificadores. Posteriormente, el diagrama de Nemenyi permitió identificar grupos de clasificadores sin diferencias estadísticamente significativas (conectados por líneas horizontales). Se observa que el método propuesto (GNLPDA, rango 4.05) se sitúa en un clúster intermedio; si bien comparte grupo estadístico con modelos de SVM polinomiales y Gaussianos, no presenta diferencias significativas frente a los métodos lineales (LDA/QDA) de mejor rango global. No obstante, destaca la superioridad absoluta de GNLPDA sobre el método KFDA ($p<0.05$), consolidando su posición como una alternativa competitiva y estadísticamente robusta frente al estado del arte en este entorno de baja complejidad.


\begin{figure}[htbp]
	\centering
	\includegraphics[width=0.8\textwidth]{CDplot_Nemenyi_1.2.penguins_long_graficos_campeones.png}
	\caption{Análisis de significancia estadística post-hoc y corrección por enmascaramiento
		mediante distancias críticas de Nemenyi para el conjunto de datos Palmer penguins.}
	\label{fig:R21}
\end{figure}

Para profundizar en la evaluación del rendimiento frente a arquitecturas no lineales,
se realizó una comparativa excluyendo los métodos deterministas lineales. La prueba de
Friedman ratifica la existencia de diferencias estadísticamente significativas en este
subconjunto ($\chi^2= 31.391, p< 0.001$). El diagrama de Nemenyi resultante evidencia que
el método propuesto (GNLPDA, rango 2.77) se integra al clúster de máxima eficacia
estadística, compartiendo línea de desempeño con SVM(POL2) y SVM(Gauss). Este
resultado es particularmente valioso, pues confirma que GNLPDA logra equiparar la
potencia de clasificación de las máquinas de vectores de soporte sin sacrificar la
interpretabilidad de su espacio de características. En contraste, KFDA (rango 4.27) se aleja
significativamente de este grupo.


\subsection{Wines}

El conjunto de datos Wine fue desarrollado a partir de análisis químicos realizados sobre vinos producidos en una misma región de Italia y derivados de tres cultivares diferentes, siendo posteriormente incorporado al repositorio UCI Machine Learning Repository por Stefan Aeberhard y colaboradores \cite{coomans1992comparison}. El conjunto está conformado por 178 observaciones distribuidas en tres clases correspondientes a los distintos tipos de vino: clase 1 ($33.1\%$), clase 2 ($39.9\%$) y clase 3 ($27.0\%$). Cada observación se describe mediante trece variables continuas relacionadas con propiedades fisicoquímicas de la muestra, entre ellas el contenido de alcohol, ácido málico, cenizas, magnesio, fenoles totales, flavonoides, intensidad de color y prolina.

Desde la perspectiva del reconocimiento de patrones, el conjunto Wine constituye un referente clásico para la evaluación de métodos de clasificación y reducción de dimensionalidad. Los datos corresponden a un problema multiclase bien definido, en el que las trece variables químicas proporcionan información discriminativa suficiente para distinguir entre los tres cultivares de vino. De hecho, la documentación original del conjunto lo describe como un problema de clasificación con una estructura de clases bien comportada, lo que lo convierte en un escenario apropiado para la validación inicial de nuevos clasificadores, aunque con un nivel de complejidad moderado. Asimismo, la presencia de trece características continuas en un espacio de mayor dimensionalidad que otros conjuntos de referencia favorece el análisis de técnicas de proyección y extracción de características, permitiendo estudiar la capacidad de los métodos para preservar la estructura discriminativa de los datos en representaciones de menor dimensión. Adicionalmente, el conjunto presenta un ligero desbalance entre clases y no contiene valores perdidos, por lo que resulta útil para comparar el desempeño de distintos enfoques bajo condiciones controladas y ampliamente documentadas en la literatura.

A continuación se muestra mediante gráficos de dispersión el comportamiento de las observaciones por clase en pares de variables:


\begin{figure}[H]
	\centering
	\includegraphics[width=0.5\textwidth]{3DES.png}
	\caption{Gráficos de dispersión por pares de Wines.}
	\label{fig:D3}
\end{figure}

La figura \ref{fig:D3}  muestra el comportamiento de las clases de vinos mediante gráficos de dispersión por partes, salta a la vista que existen pares de variables donde la separabilidad entre clases es más evidente que en otras, por ejemplo, entre las variables V1 y V7 las clases se encuentran bien separadas con un grado mínimo de traslape. Mientras que al observar pares como V5 y V10, se observa como las observaciones por clase se encuentran traslapadas entre ellas. \\


A continuación se obtuvieron las proyecciones visuales mediante la reducción de características con LDA, KFDA y GNLPDA.

\begin{figure}[H]
	\centering
	
	% Primera fila
	\begin{subfigure}[b]{0.3\textwidth}
		\centering
		\includegraphics[width=\textwidth]{mejor_escenario_1.4.wines_1_mass_2d.png}
	\end{subfigure}
	\hfill
	\begin{subfigure}[b]{0.3\textwidth}
		\centering
		\includegraphics[width=\textwidth]{mejor_escenario_1.4.wines_2_kfda_2d.png}
	\end{subfigure}
	\hfill
	\begin{subfigure}[b]{0.3\textwidth}
		\centering
		\includegraphics[width=\textwidth]{mejor_escenario_1.4.wines_3_gnlpda_2d.png}
	\end{subfigure}
	
	\caption{Comparativa geométrica de los espacios de proyección para el conjunto de datos
		Wines.}
	\label{fig:R7}
\end{figure}

La Figura \ref{fig:R7} muestra los espacios latentes generados para el conjunto de datos Wines.
Se observa que la aproximación lineal (LDA MASS) logra una separación geométrica
exitosa, preservando la varianza natural de los datos y alcanzando una exactitud perfecta. En
contraste, el mapeo no lineal implícito (KFDA) exhibe una ligera degradación en su
rendimiento predictivo, lo cual sugiere que la compresión extrema de la dispersión intra-clase
inducida por el kernel gaussiano genera un leve sobreajuste en las fronteras de decisión.
Resulta crítico destacar el comportamiento del método propuesto (GNLPDA), el cual
recupera la exactitud perfecta mediante la selección exclusiva de 5 términos de interacción
polinomial. Este hallazgo demuestra la capacidad del algoritmo para superar la robustez
predictiva de enfoques complejos, aislando interacciones de orden superior.

\begin{figure}[htbp]
	\centering
	\includegraphics[width=0.8\textwidth]{Boxplot_1.4.wines_Accuracy.png}
	\caption{Distribución empírica del rendimiento predictivo para el conjunto de datos Wines.}
	\label{fig:R22}
\end{figure}


La Figura \ref{fig:R22} sintetiza el desempeño de los clasificadores, se observa que los modelos
SVM y los métodos lineales (LDA/QDA) alcanzan una estabilización superior, con rangos
intercuartílicos compactos y una exactitud que supera el 0.99. Por su parte, el método
propuesto (GNLPDA) exhibe una mayor dispersión en su rendimiento, con un rango
intercuartílico desplazado por debajo de las arquitecturas SVM y los métodos clasicos y la
presencia de una varianza predictiva más pronunciada. Pese a esta brecha, el modelo sostiene
una capacidad discriminante superior a métodos como KFDA, validando su competitividad
frente a las aproximaciones basadas en kernel, aun cuando la parsimonia en la selección de
variables incrementa la exigencia sobre el proceso evolutivo.


\begin{figure}[htbp]
	\centering
	\includegraphics[width=0.8\textwidth]{CDplot_Nemenyi_1.4.wines_long_graficos_todos.png}
	\caption{Análisis de significancia estadística post-hoc mediante distancias críticas de
		Nemenyi para el conjunto de datos Wines.}
	\label{fig:R23}
\end{figure}

En este escenario de alta dimensionalidad, el test de Friedman reporta un estadístico
$\chi^2=94.416 (p< 0.001)$, lo que confirma la existencia de diferencias significativas en el
rendimiento de los algoritmos. El diagrama de Nemenyi resultante ilustra una clara
estratificación: mientras que los métodos lineales (LDA/QDA) y las SVM de configuración
polinomial y gaussiana ocupan los rangos superiores, el método GNLPDA se posiciona en
un nivel de desempeño intermedio, diferenciándose estadísticamente de los clasificadores
con menor rango. Si bien el modelo no alcanza la competitividad de los métodos de kernel
altamente optimizados en este dataset específico, su rango global es superior al de KFDA y
en líneas generales, se mantiene competitivo.

\begin{figure}[htbp]
	\centering
	\includegraphics[width=0.8\textwidth]{CDplot_Nemenyi_1.4.wines_long_graficos_campeones.png}
	\caption{Análisis de significancia estadística post-hoc y corrección por enmascaramiento
		mediante distancias críticas de Nemenyi para el conjunto de datos Wines.}
	\label{fig:R24}
\end{figure}

Con el propósito de aislar la capacidad discriminante de los clasificadores no lineales,
se realizó una comparativa excluyendo los métodos lineales. La prueba de Friedman ratifica
la existencia de diferencias estadísticamente significativas ( $\chi^2 = 68.914, p<0.001)$.
El diagrama de Nemenyi resultante revela una jerarquía clara: los modelos SVM (con kernels polinomiales y gaussianos) lideran el rendimiento estadístico, manteniéndose en el clúster de máxima eficacia. En este escenario, el método propuesto (GNLPDA, rango 3.82) se sitúa en una posición intermedia, mostrando una brecha significativa respecto a las SVM, aunque conservando una ventaja estadística determinante frente al modelo KFDA (rango 4.50). 



\subsection{Cancer Wisconsin}

El conjunto de datos Wisconsin Breast Cancer (Diagnostic) fue desarrollado por W. H. Wolberg, W. N. Street y O. L. Mangasarian como parte de investigaciones orientadas al diagnóstico asistido de cáncer de mama, y posteriormente fue descrito en la literatura por O. L. Mangasarian y colaboradores \cite{mangasarian1995breast}. El conjunto está conformado por 569 observaciones correspondientes a muestras de tejido mamario obtenidas mediante aspiración con aguja fina. Las observaciones se distribuyen en dos clases diagnósticas: tumores benignos ($62.7\%$) y tumores malignos ($37.3\%$). Cada muestra se representa mediante 30 variables continuas calculadas a partir de imágenes digitalizadas de los núcleos celulares presentes en el tejido, incluyendo medidas relacionadas con el radio, textura, perímetro, área, suavidad, compacidad, concavidad, simetría y dimensión fractal.

Desde la perspectiva del reconocimiento de patrones, el conjunto Wisconsin Breast Cancer constituye un referente ampliamente utilizado para la evaluación de métodos de clasificación supervisada y reducción de dimensionalidad. Las características fueron diseñadas para capturar propiedades morfológicas relevantes de los núcleos celulares, proporcionando una representación numérica con alto contenido discriminativo entre las clases benigna y maligna. Asimismo, la dimensionalidad relativamente elevada del conjunto, dada por sus treinta variables continuas, favorece el estudio de técnicas de selección de características, extracción de variables y proyección a espacios de menor dimensión. Adicionalmente, la distribución no uniforme entre clases introduce un grado moderado de desbalance que permite analizar la sensibilidad de los métodos ante distribuciones de datos asimétricas. En conjunto, estas características hacen de Wisconsin Breast Cancer un caso de estudio apropiado para comparar la capacidad discriminativa de distintos métodos de clasificación y evaluar la calidad de las representaciones obtenidas mediante técnicas de reducción de dimensionalidad.

A continuación se muestra mediante gráficos de dispersión el comportamiento de las observaciones por clase en pares de variables:


\begin{figure}[H]
	\centering
	\includegraphics[width=0.5\textwidth]{7DES.png}
	\caption{Gráficos de dispersión por pares de Cancer Wisconsin.}
	\label{fig:D7}
\end{figure}

En la figura \ref{fig:D7} se observa el comportamiento de las clases por pares de variables. 
Destaca el comportamiento de traslape que se puede observar en pares de variables como $smoothness2$ y $symetmetry2$ donde las observaciones están traslapadas. Así como tambien existen pares de variables que muestran una separación más clara de las clases como es el caso del gráfico de dispersión entre $perimeter2$ y $radius3$.


A continuación se obtuvieron las proyecciones visuales mediante la reducción de características con LDA, KFDA y GNLPDA.

\begin{figure}[H]
	\centering
	
	% Primera fila
	\begin{subfigure}[b]{0.3\textwidth}
		\centering
		\includegraphics[width=\textwidth]{mejor_escenario_4.1.CancerWisconsin_1_mass_2d.png}
	\end{subfigure}
	\hfill
	\begin{subfigure}[b]{0.3\textwidth}
		\centering
		\includegraphics[width=\textwidth]{mejor_escenario_4.1.CancerWisconsin_2_kfda_2d.png}
	\end{subfigure}
	\hfill
	\begin{subfigure}[b]{0.3\textwidth}
		\centering
		\includegraphics[width=\textwidth]{mejor_escenario_4.1.CancerWisconsin_3_gnlpda_2d.png}
	\end{subfigure}
	
	\caption{Comparativa geométrica de los espacios de proyección para el conjunto de datos
		Cancer Wisconsin.}
	\label{fig:R8}
\end{figure}

La Figura \ref{fig:R8} presenta las distribuciones de densidad estimadas para el escenario de
clasificación binaria en el conjunto de datos de Cancer Wisconsin. Se observa que el enfoque
lineal (LDA MASS) alcanza un alto rendimiento predictivo (0.982), manteniendo una
transición suave entre las clases con un área de solapamiento natural. Por su parte, el KFDA
empata esta exactitud, pero su espacio latente exhibe una severa deformación multimodal
con picos de densidad muy marcados; este comportamiento sugiere que pese a lograr un buen
desempeño, el método deforma el espacio de una manera muy agresiva. En contraste, el
método propuesto (GNLPDA) iguala la exactitud óptima de 0.982 restaurando una
distribución de densidad unimodal, fluida, suave y estadísticamente coherente. Esta
separación geométrica robusta se logra mediante la selección explícita de 16 términos
polinomiales, demostrando que GNLPDA evade las distorsiones geométricas. de los mapeos
implícitos sin sacrificar capacidad predictiva.


\begin{figure}[htbp]
	\centering
	\includegraphics[width=0.8\textwidth]{Boxplot_4.1.CancerWisconsin_Accuracy.png}
	\caption{Distribución empírica del rendimiento predictivo para el conjunto de datos Cancer
		Wisconsin.}
	\label{fig:R25}
\end{figure}

La Figura \ref{fig:R25} resume el desempeño de los clasificadores en un problema de
diagnóstico binario. Se observa que el método propuesto (GNLPDA) exhibe un rendimiento
sobresaliente, superando consistentemente en mediana a los enfoques lineales (LDA/QDA)
y consolidándose por encima de la inestabilidad operativa mostrada por el modelo KFDA. Si
bien SVM-POL2 alcanza una mediana ligeramente superior, GNLPDA mantiene un rango intercuartílico notablemente compacto y una dispersión acotada, lo que refleja una alta resiliencia y predictibilidad ante las variaciones en las muestras de entrenamiento.

\begin{figure}[H]
	\centering
	\includegraphics[width=0.8\textwidth]{CDplot_Nemenyi_4.1.CancerWisconsin_long_graficos_todos.png}
	\caption{ Análisis de significancia estadística post-hoc mediante distancias críticas de
		Nemenyi para el conjunto de datos Cancer Wisconsin.}
	\label{fig:R26}
\end{figure}

El test de Friedman reporta un estadístico $\chi^2 = 64.27 (p<0.001)$ que confirma diferencias estadísticamente significativas entre el desempeño de los clasificadores. El diagrama de Nemenyi evidencia una estructura de rangos donde los modelos SVM(POL2) y SVM(POL3) ocupan la posición de mayor eficacia estadística. El método propuesto
(GNLPDA, rango 3.70) se ubica en un clúster intermedio, mostrando una competitividad
estadística superior frente a los modelos LDA y KFDA, y manteniéndose a una distancia
crítica de los modelos líderes. Este resultado valida la robustez de la aproximación
estocástica, la cual logra un posicionamiento competitivo frente a la mayoría de los métodos
evaluados sin incurrir en la fragilidad paramétrica de las aproximaciones de kernel
tradicional.

\begin{figure}[H]
	\centering
	\includegraphics[width=0.8\textwidth]{CDplot_Nemenyi_4.1.CancerWisconsin_long_graficos_campeones.png}
	\caption{Análisis de significancia estadística post-hoc y corrección por enmascaramiento
		mediante distancias críticas de Nemenyi para el conjunto de datos Cancer Wisconsin.}
	\label{fig:R27}
\end{figure}

Al evaluar exclusivamente el subconjunto de métodos no lineales, el test de Friedman
($\chi^2= 44.898 ,p<0.001 $) ratifica la existencia de diferencias estadísticamente
significativas. El diagrama de Nemenyi detalla la jerarquía de desempeño: las arquitecturas
SVM(POL2) y SVM(POL3) mantienen el rango superior, mientras que el método propuesto
(GNLPDA, rango 3.13) se posiciona en un nivel intermedio. La relevancia de este resultado
reside en que, a pesar de la brecha estadística frente a los modelos polinomiales de mejor
rango, GNLPDA demuestra un comportamiento estadísticamente indiferenciado respecto al
grupo de alto rendimiento.


\subsection{Ionosfera}

El conjunto de datos \textit{Ionosphere} fue originalmente introducido por V. G. Sigillito y colaboradores en 1989 \cite{sigillito1989classification}. El conjunto consiste en 351 observaciones correspondientes a retornos de radar de alta frecuencia dirigidos a electrones libres en la ionosfera, distribuidas en dos clases: retornos estructurados o de interés (\textit{good}, $64.1\%$) y retornos no estructurados (\textit{bad}, $35.9\%$). Cada observación consta de 34 variables numéricas continuas que representan las partes real e imaginaria provenientes de una función de autocorrelación aplicada a patrones de multipulso de las señales electromagnéticas (comúnmente se descarta una variable constante con valor cero en la primera componente continua, resultando en 33 predictores efectivos).  Desde la perspectiva del reconocimiento de patrones, el conjunto Ionosphere presenta propiedades relevantes para la evaluación experimental de clasificadores y algoritmos de selección o reducción de dimensionalidad en espacios de moderada a alta dimensión. A diferencia de conjuntos con límites de decisión simples, sus características exhiben patrones complejos y no lineales con un nivel apreciable de solapamiento interclase en el espacio proyectado, lo que demanda métodos capaces de capturar interacciones de mayor orden para una separación efectiva. Adicionalmente, la desproporción entre las clases introduce un escenario de desbalance moderado que permite examinar la capacidad de generalización y la robustez de las técnicas analizadas ante representaciones asimétricas del espacio de características. En conjunto, estas propiedades convierten a \textit{Ionosphere} en un caso de estudio idóneo para evaluar el rendimiento de modelos de clasificación y la preservación de estructura discriminativa mediante técnicas de proyección continua.

A continuación se muestra mediante gráficos de dispersión el comportamiento de las observaciones por clase en pares de variables:

\begin{figure}[H]
	\centering
	\includegraphics[width=0.5\textwidth]{8DES.png}
	\caption{Gráficos de dispersión por pares de Ionosphere.}
	\label{fig:D8}
\end{figure}

El la figura \ref{fig:D8} se puede observar como para todos los pares de variables las observaciones se encuentran ampliamente dispersar y sin mostrar indicios de una separación lineal definida.


A continuación se obtuvieron las proyecciones visuales mediante la reducción de características con LDA, KFDA y GNLPDA.


\begin{figure}[htbp]
	\centering
	
	% Primera fila
	\begin{subfigure}[b]{0.3\textwidth}
		\centering
		\includegraphics[width=\textwidth]{mejor_escenario_4.2.Ionosphere_1_mass_2d.png}
	\end{subfigure}
	\hfill
	\begin{subfigure}[b]{0.3\textwidth}
		\centering
		\includegraphics[width=\textwidth]{mejor_escenario_4.2.Ionosphere_2_kfda_2d.png}
	\end{subfigure}
	\hfill
	\begin{subfigure}[b]{0.3\textwidth}
		\centering
		\includegraphics[width=\textwidth]{mejor_escenario_4.2.Ionosphere_3_gnlpda_2d.png}
	\end{subfigure}
	
	\caption{Comparativa geométrica de los espacios de proyección para el conjunto de datos
		Ionosphere.}
	\label{fig:R9}
\end{figure}

La Figura \ref{fig:R9} ilustra las distribuciones de densidad para un escenario caracterizado
por una alta complejidad topológica y un severo solapamiento de clases. Se advierte que el
enfoque lineal (LDA MASS) resulta bueno pero insuficiente para capturar la frontera de
decisión, exhibiendo una extensa área de ambigüedad predictiva (0.904). Por su parte, el
mapeo no lineal (KFDA) incrementa la exactitud empírica (0.962), pero a expensas de inducir
una extrema deformación multimodal en el espacio latente; los agudos picos de densidad
evidencian que el kernel gaussiano está sobreajustando la geometría para aislar muestras
específicas. En agudo contraste, el método propuesto (GNLPDA) mantiene un rendimiento
altamente competitivo (0.952) penalizando estas distorsiones y restaurando distribuciones de
densidad fluidas y generalizables. Mediante la selección explícita de 21 términos
polinomiales, GNLPDA demuestra su capacidad para lidiar con el solapamiento severo
construyendo un margen de separación matemáticamente interpretable.

\begin{figure}[htbp]
	\centering
	\includegraphics[width=0.8\textwidth]{Boxplot_4.2.Ionosphere_Accuracy.png}
	\caption{Distribución empírica del rendimiento predictivo para el conjunto de datos
		Ionosphere.}
	\label{fig:R28}
\end{figure}


La Figura \ref{fig:R28} resume el desempeño de los clasificadores en un entorno caracterizado por una alta dimensionalidad de características. El método propuesto (GNLPDA) exhibe una estabilidad destacable, con un rango intercuartílico comparable al de los métodos basados en kernel. Es preciso señalar que, en este escenario, el modelo QDA resultó computacionalmente inviable, reportando una exactitud nula debido a la singularidad o falta de rango completo en la matriz de covarianza requerida para la discriminación cuadrática. Frente a esta limitación del modelo determinista, GNLPDA demuestra una robustez operativa superior ante escenarios complicados y de alta dimensionalidad.

\begin{figure}[htbp]
	\centering
	\includegraphics[width=0.8\textwidth]{CDplot_Nemenyi_4.2.Ionosphere_long_graficos_todos.png}
	\caption{Análisis de significancia estadística post-hoc mediante distancias críticas de
		Nemenyi para el conjunto de datos Ionosphere.}
	\label{fig:R29}
\end{figure}

Tras la exclusión del modelo QDA por inviabilidad computacional, el análisis inferencial subraya una estructura de rendimiento heterogénea ($\chi^2=88.363 ,p<0.001$). El diagrama de Nemenyi muestra que, mientras los modelos SVM(Gauss), KFDA y SVM(POL2) se posicionan en el clúster de mayor eficacia estadística, el método propuesto (GNLPDA, rango 3.95) ocupa una posición intermedia. Es notable que, a pesar de la complejidad del dominio, el método logra una separación estadística clara frente a los modelos de menor rango, consolidándose como una alternativa funcional y robusta allí donde otros clasificadores estadísticos tradicionales enfrentan fallos de convergencia numérica. 

\begin{figure}[htbp]
	\centering
	\includegraphics[width=0.8\textwidth]{CDplot_Nemenyi_4.2.Ionosphere_long_graficos_campeones.png}
	\caption{Análisis de significancia estadística post-hoc y corrección por enmascaramiento
		mediante distancias críticas de Nemenyi para el conjunto de datos Ionosphere.}
	\label{fig:R30}
\end{figure}

Al restringir el análisis exclusivamente a los clasificadores no lineales, el test de
Friedman corrobora la persistencia de diferencias estadísticamente 
($\chi^2 = 88.363,p<0.001$). El diagrama de Nemenyi resultante revela una jerarquía donde las configuraciones de SVM (Gauss y POL2) y KFDA dominan la eficacia estadística. El método propuesto (GNLPDA) mantiene su posición intermedia, evidenciando que, aun prescindiendo de los métodos lineales, la brecha de rendimiento frente a las máquinas de soporte vectorial permanece estable.


\subsection{Cebolla 3D}

El conjunto de datos sintético Cebolla3D consiste en 300 observaciones generadas sintéticamente y distribuidas equitativamente en tres clases concéntricas en un espacio tridimensional ($p=3$). Cada clase corresponde a un cascarón esférico definido por un radio característico perturbado mediante un factor de ruido gaussiano. Desde la perspectiva del reconocimiento de patrones, constituye un caso de prueba geométrico controlado de baja dimensión pero estrictamente no lineal, diseñado para evaluar la capacidad de proyección y clasificación ante manifolds cerrados donde las fronteras de decisión lineales fracasan por completo.

A continuación se muestra mediante gráficos de dispersión el comportamiento de las observaciones por clase en pares de variables:

\begin{figure}[H]
	\centering
	\includegraphics[width=0.5\textwidth]{4DES.png}
	\caption{Gráficos de dispersión por pares de Cebolla3D.}
	\label{fig:D4}
\end{figure}

El la figura \ref{fig:D4} se observa esta formación característica de capas concéntricas con diferentes grados de dispersión que no son linealmente separables.


A continuación se obtuvieron las proyecciones visuales mediante la reducción de características con LDA, KFDA y GNLPDA.

\begin{figure}[htbp]
	\centering
	
	% Primera fila
	\begin{subfigure}[b]{0.3\textwidth}
		\centering
		\includegraphics[width=\textwidth]{mejor_escenario_3.1.Cebolla3D_1_mass_2d.png}
	\end{subfigure}
	\hfill
	\begin{subfigure}[b]{0.3\textwidth}
		\centering
		\includegraphics[width=\textwidth]{mejor_escenario_3.1.Cebolla3D_2_kfda_2d.png}
	\end{subfigure}
	\hfill
	\begin{subfigure}[b]{0.3\textwidth}
		\centering
		\includegraphics[width=\textwidth]{mejor_escenario_3.1.Cebolla3D_3_gnlpda_2d.png}
	\end{subfigure}
	
	\caption{Comparativa geométrica de los espacios de proyección para el conjunto de datos
		cebolla 3D.}
	\label{fig:R10}
\end{figure}

La Figura \ref{fig:R10} evidencia el comportamiento algorítmico frente a la búsqueda de una
frontera de separación no lineal. Como es esperable ante estructuras de distribución
concéntrica, el enfoque estrictamente lineal (LDA MASS) fracasa en la proyección del
espacio latente separable, colapsando las clases en un solapamiento masivo que resulta en
una exactitud predictiva deficiente (0.433). En contraparte, el mapeo implícito (KFDA) logra
desdoblar exitosamente esta complejidad geométrica utilizando el kernel gaussiano,
alcanzando una separación casi perfecta (0.989). El hallazgo analítico más trascendental se
observa en el método propuesto (GNLPDA), el cual iguala el rendimiento máximo del kernel
mediante la selección de únicamente 4 términos. Es particularmente destacable mencionar
que el algoritmo seleccionó los predictores cuadráticos; este comportamiento demuestra que
la búsqueda estocástica identifica con precisión las transformaciones polinomiales necesarias
para linealizar la estructura concéntrica, logrando proyectar fronteras de decisión óptimas
dentro de un espacio de características explícito, parsimonioso y totalmente interpretable.

\begin{figure}[htbp]
	\centering
	\includegraphics[width=0.8\textwidth]{Boxplot_3.1.Cebolla3D_Accuracy.png}
	\caption{Distribución empírica del rendimiento predictivo para el conjunto de datos
		cebolla 3D.}
	\label{fig:R31}
\end{figure}

La Figura \ref{fig:R31} ilustra el desempeño de los clasificadores ante una geometría no lineal
de tres dimensiones. Se observa un colapso operativo del clasificador lineal (LDA), cuya
exactitud basada en la mediana se desploma por debajo de 0.40, validando la inoperancia de
los enfoques lineales ante distribuciones de clases no linealmente separables. En contraste,
el método propuesto (GNLPDA) exhibe una alta eficacia predictiva, alcanzando medianas superiores a 0.90 y manteniéndose competitivo frente a los modelos basados en kernel y máquinas de vectores de soporte. Es destacable que QDA logra obtener desempeños competitivos ante este escenario, sosteniendo un rango intercuartílico compacto y una mediana superior a 0.8.

\begin{figure}[H]
	\centering
	\includegraphics[width=0.8\textwidth]{CDplot_Nemenyi_3.1.Cebolla3D_long_graficos_todos.png}
	\caption{Análisis de significancia estadística post-hoc mediante distancias críticas de
		Nemenyi para el conjunto de datos cebolla 3D.}
	\label{fig:R32}
\end{figure}

La prueba de Friedman confirma la existencia de diferencias estadísticamente significativas en el rendimiento de los algoritmos ($\chi^2=142.551,p<0.001 $). El diagrama de Nemenyi revela una segmentación clara: las configuraciones de SVM (POL2, Gauss y POL3) integran el clúster de máxima eficacia estadística. El método propuesto (GNLPDA, rango 4.27) se posiciona en un nivel intermedio, mostrando un desempeño estadísticamente indiferenciado respecto al KFDA, siendo que en ultima instancia, se genera un grupo donde están LDA y QDA, siendo los de menor rendimiento en comparación.

\begin{figure}[H]
	\centering
	\includegraphics[width=0.8\textwidth]{CDplot_Nemenyi_3.1.Cebolla3D_long_graficos_campeones.png}
	\caption{Análisis de significancia estadística post-hoc y corrección por enmascaramiento
		mediante distancias críticas de Nemenyi para el conjunto de datos Cebolla 3D.}
	\label{fig:R33}
\end{figure}


Al evaluar exclusivamente el subconjunto de métodos no lineales, la prueba de
Friedman ratifica la existencia de diferencias estadísticamente significativas
($\chi^2 = 89.354 ,p<0.001$). El diagrama de Nemenyi detalla una clara jerarquía:
las máquinas de vectores de soporte (SVM-POL2, SVM-Gauss y SVM-POL3) integran el
clúster de máxima eficacia estadística. El método propuesto (GNLPDA, rango 4.08) se
posiciona en un nivel intermedio, manteniendo una equivalencia estadística frente al KFDA
(rango 4.23). 

\subsection{Cebolla 10D}

El conjunto de datos sintético Cebolla10D consta de 300 observaciones pertenecientes a tres clases concéntricas proyectadas en un espacio de 10 dimensiones ($p=10$). El conjunto conserva la estructura de hiper-esferas imbricadas con perturbación gaussiana radial. Desde la perspectiva del reconocimiento de patrones, este conjunto escala la dificultad geométrica al añadir dimensiones ruidosas u ortogonales al espacio de la variedad, permitiendo examinar la estabilidad y robustez de los algoritmos de selección de características y transformación no lineal a medida que aumenta la dimensionalidad del espacio original.

A continuación se muestra mediante gráficos de dispersión el comportamiento de las observaciones por clase en pares de variables:

\begin{figure}[H]
	\centering
	\includegraphics[width=0.5\textwidth]{5DES.png}
	\caption{Gráficos de dispersión por pares de Cebolla10D.}
	\label{fig:D5}
\end{figure}

El la figura \ref{fig:D5} se observa la dispersión presente por cada clase de forma que los traslapes hacen que no se observe una separación clara entre las clases.


A continuación se obtuvieron las proyecciones visuales mediante la reducción de características con LDA, KFDA y GNLPDA.


\begin{figure}[htbp]
	\centering
	
	% Primera fila
	\begin{subfigure}[b]{0.3\textwidth}
		\centering
		\includegraphics[width=\textwidth]{mejor_escenario_3.2.Cebolla10D_1_mass_2d.png}
	\end{subfigure}
	\hfill
	\begin{subfigure}[b]{0.3\textwidth}
		\centering
		\includegraphics[width=\textwidth]{mejor_escenario_3.2.Cebolla10D_2_kfda_2d.png}
	\end{subfigure}
	\hfill
	\begin{subfigure}[b]{0.3\textwidth}
		\centering
		\includegraphics[width=\textwidth]{mejor_escenario_3.2.Cebolla10D_3_gnlpda_2d.png}
	\end{subfigure}
	
	\caption{Comparativa geométrica de los espacios de proyección para el conjunto de datos
		cebolla 10D.}
	\label{fig:R11}
\end{figure}

La Figura \ref{fig:R11} ilustra el comportamiento de los clasificadores ante un escenario que
escala la complejidad de la separación concéntrica a diez dimensiones. Como resulta
previsible, la restricción del enfoque estrictamente lineal (LDA MASS) provoca un colapso
geométrico absoluto, evidenciado por un solapamiento masivo de las clases y una exactitud
predictiva apenas superior al azar (0.511). En contraste, la expansión implícita del kernel
gaussiano (KFDA) logra sortear esta geometría no lineal, alcanzando una separación perfecta
(1.0). No obstante, se observa en la proyección del método propuesto que sostiene un
rendimiento casi óptimo (0.989) a través de un subconjunto explícito de 24 términos. Resulta
excepcionalmente destacable que la búsqueda estocástica logró identificar e incorporar la
totalidad de los predictores cuadráticos de las variables originales. Esta evidencia empírica
demuestra de forma contundente que el algoritmo es capaz de reconstruir analíticamente la
ecuación de una hiperesfera en espacios de alta dimensionalidad, logrando linealizar el
problema mediante una base polinomial exacta y totalmente transparente.

\begin{figure}[htbp]
	\centering
	\includegraphics[width=0.8\textwidth]{Boxplot_3.2.Cebolla10D_Accuracy.png}
	\caption{Distribución empírica del rendimiento predictivo para el conjunto de datos Cancer
		cebolla 10D.}
	\label{fig:R34}
\end{figure}


La Figura \ref{fig:R34} muestra el desempeño de los clasificadores ante una geometría no
lienalmente separable de diez dimensiones. El colapso del modelo lineal (LDA) es absoluto,
con una exactitud mediana inferior a 0.45. Mientras que los modelos de kernel (KFDA y
SVM-Gauss) alcanzan los niveles máximos de eficacia predictiva, el método propuesto
(GNLPDA) mantiene un rendimiento sólido, con una mediana superior a 0.88. Cabe destacar
que el modelo QDA se mantiene con una mediana superior al 0.8, donde si bien, no figura
como los mejores, se mantiene competente.

\begin{figure}[H]
	\centering
	\includegraphics[width=0.8\textwidth]{CDplot_Nemenyi_3.2.Cebolla10D_long_graficos_todos.png}
	\caption{Análisis de significancia estadística post-hoc mediante distancias críticas de
		Nemenyi para el conjunto de datos cebolla 10D.}
	\label{fig:R35}
\end{figure}

La prueba de Friedman confirma la existencia de diferencias estadísticamente
significativas en el rendimiento de los algoritmos ($\chi^2= 144.857 ,p<0.001$).
El diagrama de Nemenyi revela una estratificación jerárquica: el clúster de mayor
eficacia estadística está integrado por los modelos KFDA y SVM (Gauss). El método
propuesto (GNLPDA, rango 4.33) se posiciona en un nivel intermedio, mostrando
equivalencia estadística frente a los modelos SVM(POL2) y SVM(POL3). Finalmente, los
modelos QDA y LDA ocupan los rangos inferiores, reafirmando la superioridad de los
enfoques no lineales frente a estructuras de discriminación lineal o cuadrática simple en este
entorno de alta dimensionalidad.

\begin{figure}[H]
	\centering
	\includegraphics[width=0.8\textwidth]{CDplot_Nemenyi_3.2.Cebolla10D_long_graficos_campeones.png}
	\caption{Análisis de significancia estadística post-hoc y corrección por enmascaramiento
		mediante distancias críticas de Nemenyi para el conjunto de datos cebolla 10D.}
	\label{fig:R36}
\end{figure}

Al evaluar exclusivamente el subconjunto de métodos no lineales, el test de Friedman
corrobora la existencia de diferencias estadísticamente significativas ($\chi^2 = = 89.354,p<0.001$). El diagrama de Nemenyi muestra una separación clara: el clúster de mayor eficacia estadística está compuesto por KFDA y SVM(Gauss). El
método propuesto (GNLPDA, rango 3.93) se posiciona junto a las variantes polinomiales de SVM (POL2 y POL3) en un nivel de desempeño intermedio. 

\subsection{Cebolla 20D}

El conjunto de datos sintético Cebolla20D comprende 300 observaciones organizadas en tres clases concéntricas dentro de un espacio de 20 dimensiones ($p=20$). Al igual que las variantes anteriores, el proceso de generación emplea hiper-esferas con ruido estocástico integrado. Desde la perspectiva del reconocimiento de patrones, representa un escenario de prueba en moderada-alta dimensión donde la dispersión del espacio geométrico (maldición de la dimensionalidad) complica la identificación de manifolds concéntricos, resultando idóneo para evaluar la preservación de la estructura discriminativa y la capacidad de generalización de los clasificadores no lineales.


A continuación se muestra mediante gráficos de dispersión el comportamiento de las observaciones por clase en pares de variables:

\begin{figure}[H]
	\centering
	\includegraphics[width=0.5\textwidth]{6DES.png}
	\caption{Gráficos de dispersión por pares de Cebolla20D.}
	\label{fig:D6}
\end{figure}

El la figura \ref{fig:D6} se observa un fuerte traslape entre las clases así como un alto grado de dispersión que no permite encontrar limites de separación claros entre una clase y otra.

A continuación se obtuvieron las proyecciones visuales mediante la reducción de características con LDA, KFDA y GNLPDA.

\begin{figure}[htbp]
	\centering
	
	% Primera fila
	\begin{subfigure}[b]{0.3\textwidth}
		\centering
		\includegraphics[width=\textwidth]{mejor_escenario_3.3.Cebolla20D_1_mass_2d.png}
	\end{subfigure}
	\hfill
	\begin{subfigure}[b]{0.3\textwidth}
		\centering
		\includegraphics[width=\textwidth]{mejor_escenario_3.3.Cebolla20D_2_kfda_2d.png}
	\end{subfigure}
	\hfill
	\begin{subfigure}[b]{0.3\textwidth}
		\centering
		\includegraphics[width=\textwidth]{mejor_escenario_3.3.Cebolla20D_3_gnlpda_2d.png}
	\end{subfigure}
	
	\caption{Comparativa geométrica de los espacios de proyección para el conjunto de datos
		cebolla 20D.}
	\label{fig:R12}
\end{figure}

La Figura \ref{fig:R12} exhibe el límite de estrés dimensional al proyectar una geometría
concéntrica de 20 variables. El clasificador lineal (LDA MASS) confirma su inoperatividad
ante esta estructura, registrando un severo solapamiento inter-clase y una exactitud deficiente
(0.600). Por el contrario, el mapeo implícito (KFDA) absorbe la complejidad proyectando
una separación perfecta (1.000). El análisis de la proyección generada por el método
propuesto (GNLPDA) revela un comportamiento de alto valor heurístico: a pesar de
enfrentarse a una explosión combinatoria en el espacio de características expandido (230
términos posibles), el algoritmo sostiene un rendimiento predictivo sobresaliente (0.933)
empleando un subconjunto de apenas 38 variables. Resulta notable que la búsqueda
estocástica logró identificar 11 de los 20 predictores cuadráticos fundamentales. Esta
evidencia empírica demuestra que, ante un escalamiento dimensional masivo, el método
prioriza un equilibrio óptimo entre la parsimonia del modelo, la extracción de los
componentes no lineales subyacentes y una alta capacidad de separación, conservando la
explicabilidad matemática frente al enfoque de caja negra del kernel.

\begin{figure}[htbp]
	\centering
	\includegraphics[width=0.8\textwidth]{Boxplot_3.2.Cebolla10D_Accuracy.png}
	\caption{Distribución empírica del rendimiento predictivo para el conjunto de datos
		cebolla 20D.}
	\label{fig:R37}
\end{figure}

La Figura \ref{fig:R37} sintetiza el desempeño de los clasificadores ante una geometría no lineal de veinte dimensiones. En este entorno, se observa una degradación significativa del rendimiento general, donde incluso las arquitecturas de kernel más robustas experimentan una mayor dispersión. El modelo lineal (LDA) confirma su ineficacia estructural, situándose en rendimientos cercanos al azar. Por su parte, el método propuesto (GNLPDA) mantiene una mediana de rendimiento competitiva, aunque con una varianza creciente producto del aumento exponencial en el espacio de búsqueda. Los modelos SVM (Gauss) y KFDA se destacan del resto por mantener los desempeños más altos del gráfico, así como un rango intercuartílico pequeño en comparación con los demás competidores.

\begin{figure}[H]
	\centering
	\includegraphics[width=0.8\textwidth]{CDplot_Nemenyi_3.3.Cebolla20D_long_graficos_todos.png}
	\caption{Análisis de significancia estadística post-hoc mediante distancias críticas de
		Nemenyi para el conjunto de datos cebolla 20D.}
	\label{fig:R38}
\end{figure}

La prueba de Friedman ($\chi^2=2 = 167.214 ,p<0.001$) corrobora la existencia de diferencias estadísticamente significativas entre los clasificadores ante la complejidad extrema del espacio de 20 dimensiones. El diagrama de Nemenyi evidencia una jerarquía donde los modelos SVM(Gauss) y KFDA lideran la eficacia estadística. El método propuesto (GNLPDA, rango 4.13) se sitúa en un nivel intermedio, manteniendo una equivalencia estadística frente al modelo QDA y SVM(POL2).

\begin{figure}[H]
	\centering
	\includegraphics[width=0.8\textwidth]{CDplot_Nemenyi_3.3.Cebolla20D_long_graficos_campeones.png}
	\caption{Análisis de significancia estadística post-hoc y corrección por enmascaramiento
		mediante distancias críticas de Nemenyi para el conjunto de datos cebolla 20D.}
	\label{fig:R39}
\end{figure}

Al evaluar exclusivamente el subconjunto de métodos no lineales, el test de Friedman corrobora la existencia de diferencias estadísticamente significativas ($\chi^2= 113.539, p<0.001$). El diagrama de Nemenyi muestra una separación clara donde los modelos SVM(Gauss) y KFDA lideran el desempeño. El método propuesto (GNLPDA, rango 3.77) se sitúa en un nivel de desempeño intermedio, manteniendo una equivalencia estadística frente a SVM(POL2). 

%%%%%%%%%%%%%%%%%%%%%%%%%%%%%%%%%%%%%%%%%%%%%%%%%%%%%%%%%%%%%%%%%%%%%%%%%%%%%%%%%%%%%%%%%%

\subsection{Pima diabetes}

El conjunto de datos Pima Indians Diabetes fue originalmente introducido por J. W. Smith y colaboradores en 1988 \cite{smith1988using}. El conjunto consiste en 768 observaciones recopiladas del Instituto Nacional de Diabetes y Enfermedades Digestivas y del Riñón, limitadas a mujeres de al menos 21 años de edad pertenecientes a la herencia indígena Pima. La variable respuesta es binaria e indica la presencia ($34.9\%$) o ausencia ($65.1\%$) de diabetes mellitus tipo 2 según los criterios del examen de la OMS. Cada observación consta de ocho variables médicas continuas y discretas: número de embarazos, concentración de glucosa en plasma, presión arterial diastólica, grosor del pliegue cutáneo del tríceps, insulina sérica de dos horas, índice de masa corporal, función de pedigrí de diabetes y edad.Desde la perspectiva del reconocimiento de patrones, el conjunto Pima Diabetes representa un escenario de prueba fundamental en el dominio biomédico para la evaluación experimental de clasificadores y métodos de reducción de dimensionalidad o selección de características. Sus predictores presentan correlaciones complejas y grados variables de no linealidad, sumado a un solapamiento considerable entre las clases en el espacio multivariado de características que dificulta la delimitación de fronteras de decisión lineales simples. Asimismo, la distribución desproporcionada de observaciones genera un escenario de desbalance moderado de clases, mientras que la presencia implícita de valores numéricamente inviables (como presiones arteriales o índices de masa corporal iguales a cero) permite examinar la sensibilidad, estabilidad y robustez de los algoritmos ante datos ruidosos e incompletos en un contexto analítico realista.

A continuación se muestra mediante gráficos de dispersión el comportamiento de las observaciones por clase en pares de variables:

\begin{figure}[H]
	\centering
	\includegraphics[width=0.5\textwidth]{2DES.png}
	\caption{Gráficos de dispersión por pares de Pima Diabetes}
	\label{fig:D2}
\end{figure}

El la figura \ref{fig:D2} como las distribuciones por clase de las variables mantienen estructuras similares, mostrando un grado de traslape y dispersión que las hace difíciles de diferenciar entre una clase y otra.


A continuación se obtuvieron las proyecciones visuales mediante la reducción de características con LDA, KFDA y GNLPDA.


\begin{figure}[htbp]
	\centering
	
	% Primera fila
	\begin{subfigure}[b]{0.3\textwidth}
		\centering
		\includegraphics[width=\textwidth]{mejor_escenario_1.3.pimadiabetes_1_mass_2d.png}
	\end{subfigure}
	\hfill
	\begin{subfigure}[b]{0.3\textwidth}
		\centering
		\includegraphics[width=\textwidth]{mejor_escenario_1.3.pimadiabetes_2_kfda_2d.png}
	\end{subfigure}
	\hfill
	\begin{subfigure}[b]{0.3\textwidth}
		\centering
		\includegraphics[width=\textwidth]{mejor_escenario_1.3.pimadiabetes_3_gnlpda_2d.png}
	\end{subfigure}
	
	\caption{Comparativa geométrica de los espacios de proyección para el conjunto de datos
		Pima diabetes.}
	\label{fig:R13}
\end{figure}

La Figura \ref{fig:R13} presenta las distribuciones de densidad estimadas para un escenario
biomédico caracterizado por un alto grado de solapamiento natural. En este contexto, el
enfoque lineal (LDA MASS) establece una línea base de rendimiento (0.817), proyectando
curvas con una dispersión amplia que refleja la varianza intrínseca de los datos fisiológicos.
Por su parte, el mapeo implícito (KFDA) presenta una leve degradación predictiva (0.809)
acompañada de un notable cambio topológico: una compresión extrema del espacio intraclase que genera picos de densidad inusualmente agudos. Esta aglomeración forzada evidencia que el kernel gaussiano tiende a colapsar las proyecciones hacia coordenadas singulares, distorsionando y perdiendo la representación de la dispersión biológica original
sin que esto se traduzca en una mejora para la separabilidad del modelo. Resulta sumamente
relevante que el método propuesto (GNLPDA) supera el rendimiento del kernel (0.813) y
evita este colapso, restaurando la fluidez y amplitud natural de las distribuciones. Esta
resiliencia geométrica se alcanza mediante la extracción de un subconjunto interpretable de
10 términos, destacando interacciones entre variables fisiológicas clave (como glucosa, edad
y masa corporal), lo que demuestra que la búsqueda estocástica modela eficazmente la
complejidad clínica sin recurrir a deformaciones extremas del espacio latente.

\begin{figure}[htbp]
	\centering
	\includegraphics[width=0.8\textwidth]{Boxplot_1.3.pimadiabetes_Accuracy.png}
	\caption{Distribución empírica del rendimiento predictivo para el conjunto de datos Pima
		diabetes.}
	\label{fig:R40}
\end{figure}

La Figura \ref{fig:R40} resume el desempeño de los clasificadores en un problema de diagnóstico biomédico con alta tasa de ruido. El método propuesto (GNLPDA) exhibe un rendimiento sobresaliente, posicionándose con la mediana de exactitud más elevada del conjunto, superando incluso a los modelos de kernel (SVM y KFDA) y lineales (LDA/QDA) aunque se observa una dispersión moderada. Por lo demás, los métodos SVM y QDA parecen mantener rendimientos similares por la posición de su mediana y su rango intercuartílico.

\begin{figure}[htbp]
	\centering
	\includegraphics[width=0.8\textwidth]{CDplot_Nemenyi_1.3.pimadiabetes_long_graficos_todos.png}
	\caption{Análisis de significancia estadística post-hoc mediante distancias críticas de
		Nemenyi para el conjunto de datos pima diabetes.}
	\label{fig:R41}
\end{figure}

La prueba de Friedman ($\chi^2=48.175 ,p<0.001$ ) ratifica la existencia de diferencias estadísticamente significativas en el desempeño de los algoritmos. El diagrama de Nemenyi muestra la formación de clústeres basados en el rendimiento estadístico: el método propuesto (GNLPDA, rango 2.58) se integra en el primer grupo junto con LDA
(rango 2.97), denotando que ambos clasificadores comparten el rango de mejor desempeño en este conjunto de datos. Por otro lado, los clasificadores SVM(POL3), SVM(Gauss), SVM(POL2) y KFDA se ubican en niveles de rango intermedios, mientras que el modelo QDA (rango 5.95) se posiciona en el extremo inferior de la jerarquía estadística. Las líneas horizontales negras y la línea verde conectan a los clasificadores que no presentan diferencias estadísticamente significativas entre sí de acuerdo con el umbral de Diferencia Crítica (CD) establecido.

\begin{figure}[htbp]
	\centering
	\includegraphics[width=0.8\textwidth]{CDplot_Nemenyi_1.3.pimadiabetes_long_graficos_campeones.png}
	\caption{Análisis de significancia estadística post-hoc y corrección por enmascaramiento
		mediante distancias críticas de Nemenyi para el conjunto de datos Pima diabetes.}
	\label{fig:R42}
\end{figure}

La prueba de Friedman ($\chi^2= 15.129 , p<0.001 $) corrobora la existencia de
diferencias estadísticamente significativas entre los clasificadores. El diagrama de Nemenyi
presenta la jerarquía de rangos observada: el método propuesto (GNLPDA, rango 2.12) se
ubica en la posición de mejor desempeño estadístico, distanciándose del resto de los
clasificadores. Los modelos SVM (POL3, rango 3.02), SVM (POL2, rango 3.12) y
SVM(Gauss, rango 3.15) conforman un grupo intermedio sin diferencias estadísticamente
significativas entre ellos. Finalmente, el clasificador KFDA (rango 3.60) se posiciona en el
extremo inferior de la jerarquía presentada. Las conexiones por líneas horizontales indican
los conjuntos de clasificadores que, bajo el umbral de Diferencia Crítica (CD), no muestran
diferencias estadísticas significativas.


\subsection{Indicadores 2015}

El conjunto de datos Indicadores 2015 fue consolidado por J. A. Navarro-Acosta y colaboradores en el 2025, a partir de mediciones sociodemográficas reales del CONEVAL y la Encuesta Intercensal de INEGI \cite{navarro2025bayesian}. El conjunto consta de 2,446 observaciones municipales organizadas en cinco clases ordinales según el Grado de Rezago Social: muy bajo, bajo, medio, alto y muy alto. Cada observación integra 12 variables continuas que abarcan 11 indicadores de carencias socioeconómicas en educación, salud, vivienda y bienes del hogar, más el índice sintético.  Desde la perspectiva del reconocimiento de patrones, representa un caso de estudio real de mediana dimensión caracterizado por un severo desbalance de clases y distribuciones con asimetría pronunciada. Estas propiedades generan solapamiento en el espacio de características, lo que permite evaluar la capacidad de los clasificadores para delimitar fronteras de decisión no lineales complejas sin asumir normalidad en los predictores.

A continuación se muestra mediante gráficos de dispersión el comportamiento de las observaciones por clase en pares de variables:

\begin{figure}[H]
	\centering
	\includegraphics[width=0.5\textwidth]{9DES.png}
	\caption{Gráficos de dispersión por pares de indicadores 2015}
	\label{fig:D9}
\end{figure}

El la figura \ref{fig:D9} se observa un grado de traslape entre las clases de forma que mantienen una baja dispersión.  

A continuación se obtuvieron las proyecciones visuales mediante la reducción de características con LDA, KFDA y GNLPDA.


\begin{figure}[H]
	\centering
	
	% Primera fila
	\begin{subfigure}[b]{0.3\textwidth}
		\centering
		\includegraphics[width=\textwidth]{mejor_escenario_5.4.indicadores2015_1_mass_2d.png}
	\end{subfigure}
	\hfill
	\begin{subfigure}[b]{0.3\textwidth}
		\centering
		\includegraphics[width=\textwidth]{mejor_escenario_5.4.indicadores2015_2_kfda_2d.png}
	\end{subfigure}
	\hfill
	\begin{subfigure}[b]{0.3\textwidth}
		\centering
		\includegraphics[width=\textwidth]{mejor_escenario_5.4.indicadores2015_3_gnlpda_2d.png}
	\end{subfigure}
	
	\caption{Comparativa geométrica de los espacios de proyección para el conjunto de datos
		Indicadores 2015 (2D).}
	\label{fig:R15}
\end{figure}

\begin{figure}[H]
	\centering
	
	% Primera fila
	\begin{subfigure}[b]{0.3\textwidth}
		\centering
		\includegraphics[width=\textwidth]{mejor_escenario_5.4.indicadores2015_1_mass_3d.png}
	\end{subfigure}
	\hfill
	\begin{subfigure}[b]{0.3\textwidth}
		\centering
		\includegraphics[width=\textwidth]{mejor_escenario_5.4.indicadores2015_2_kfda.png}
	\end{subfigure}
	\hfill
	\begin{subfigure}[b]{0.3\textwidth}
		\centering
		\includegraphics[width=\textwidth]{mejor_escenario_5.4.indicadores2015_3_gnlpda.png}
	\end{subfigure}
	
	\caption{Comparativa geométrica de los espacios de proyección para el conjunto de datos
		Indicadores 2015 (3D).}
	\label{fig:R16}
\end{figure}

Las Figuras 20 y 21 exhiben las proyecciones geométricas para un escenario
socioeconómico de cinco clases que representan niveles ordinales de rezago. Dada esta
naturaleza progresiva, la geométrica esperada es la de una trayectoria curva que transicione
de forma continua a medida que avanzan los niveles. Las vistas en 2D y 3D demuestran que
el enfoque lineal (LDA MASS) logra capturar esta progresión ordinal, estructurando dicha
curva con una exactitud destacable (0.919). En drástico contraste, el mapeo implícito
(KFDA) sufre un colapso topológico severo; el kernel gaussiano destruye por completo la
relación secuencial de los datos, provocando un solapamiento masivo que desploma el
rendimiento predictivo (0.682). El aporte analítico definitivo se consolida en el método
propuesto (GNLPDA), el cual maximiza la exactitud general (0.941) mediante la selección
de 27 términos. La proyección tridimensional evidencia que la evolución estocástica no solo
optimiza la separabilidad inter-clase, sino que refina matemáticamente esta curva, proyectando una transición geométrica fluida, parsimoniosa y perfectamente estructurada que respeta y acentúa la naturaleza ordinal del fenómeno subyacente.

\begin{figure}[H]
	\centering
	\includegraphics[width=0.8\textwidth]{Boxplot_5.4.indicadores2015_Accuracy.png}
	\caption{Distribución empírica del rendimiento predictivo para el conjunto de datos Indicadores 2015.}
	\label{fig:R43}
\end{figure}

La Figura \ref{fig:R43} detalla el comportamiento de los clasificadores en términos de exactitud
(accuracy). Los modelos SVM(POL2), SVM(POL3) y SVM(Gauss) alcanzan las medianas
más elevadas, posicionándose consistentemente por encima de 0.98. El método propuesto
(GNLPDA) exhibe una mediana superior a 0.92, mostrando una distribución concentrada con
valores de exactitud competitivos. Seguidamente, los clasificadores LDA y QDA presentan
medianas situadas en el rango de 0.87 a 0.89. Finalmente, el clasificador KFDA muestra el
menor desempeño, con una mediana de exactitud próxima a 0.65 y una mayor dispersión en
los datos en comparación con el resto de los algoritmos evaluados.

\begin{figure}[H]
	\centering
	\includegraphics[width=0.8\textwidth]{CDplot_Nemenyi_5.4.indicadores2015_long_graficos_todos.png}
	\caption{Análisis de significancia estadística post-hoc mediante distancias críticas de
		Nemenyi para el conjunto de datos Indicadores 2015.}
	\label{fig:R44}
\end{figure}

La prueba de Friedman ($\chi^2= 169.526 ,p<0.001 $) ratifica la existencia de
diferencias estadísticamente significativas en el desempeño de los algoritmos. El diagrama
de Nemenyi evidencia una jerarquía donde los modelos SVM(POL2), SVM(Gauss) y
SVM(POL3) conforman el grupo de mayor eficacia estadística. El método propuesto
(GNLPDA, rango 4.00) ocupa una posición intermedia en la jerarquía, presentando una
diferencia estadística significativa respecto al clúster líder, pero superando a los modelos
LDA, QDA y KFDA. Las conexiones mediante líneas horizontales agrupan a los
clasificadores que no presentan diferencias estadísticas significativas entre sí, de acuerdo con
la diferencia crítica (CD) calculada.

\begin{figure}[H]
	\centering
	\includegraphics[width=0.8\textwidth]{CDplot_Nemenyi_5.4.indicadores2015_long_graficos_campeones.png}
	\caption{Análisis de significancia estadística post-hoc y corrección por enmascaramiento
		mediante distancias críticas de Nemenyi para el conjunto de datos Indicadores 2015.}
	\label{fig:R45}
\end{figure}

Al restringir el análisis exclusivamente a los clasificadores no lineales, la prueba de
Friedman corrobora la existencia de diferencias estadísticamente significativas en el
desempeño de los algoritmos ($\chi^2 =103.159 , p< 0.001 $ ). El diagrama de Nemenyi detalla
una jerarquía donde los modelos SVM(POL2), SVM(Gauss) y SVM(POL3) conforman el
clúster de mayor eficacia estadística, sin diferencias significativas entre ellos. El método
propuesto (GNLPDA, rango 4.00) se posiciona en un nivel intermedio, presentando una
diferencia estadística significativa respecto al grupo líder, pero superando al clasificador
KFDA (rango 5.00), el cual se ubica en el extremo inferior de la jerarquía. Las conexiones
mediante líneas horizontales identifican a los grupos de clasificadores que, según el umbral
de Diferencia Crítica (CD), no muestran diferencias estadísticas significativas.


\subsection{Indicadores 2020}

El conjunto de datos Indicadores 2020 proviene de la misma investigación desarrollada por J. A. Navarro-Acosta y colaboradores en el 2025, incorporando mediciones reales procedentes del Censo de Población y Vivienda 2020 del INEGI \cite{navarro2025bayesian}. Este conjunto comprende 2,470 observaciones correspondientes a la totalidad de los municipios medidos en dicho periodo. Conserva la misma estructura de 12 variables continuas (11 predictores socioeconómicos y el índice sintético) distribuidas en las mismas cinco categorías del Grado de Rezago Social.  Desde la perspectiva del reconocimiento de patrones, al tratarse de un conjunto de datos real y reciente, ofrece un escenario idóneo para examinar la estabilidad y capacidad de generalización de los algoritmos ante la evolución temporal de las clases. Al compartir el mismo espacio de características que su versión de 2015 pero reflejar la transición demográfica reciente, permite evaluar la sensibilidad de las técnicas de reducción de dimensionalidad y clasificación frente a distribuciones no uniformes y esquemas de desbalance entre clases en contextos socioeconómicos actuales.

A continuación se muestra mediante gráficos de dispersión el comportamiento de las observaciones por clase en pares de variables:

\begin{figure}[H]
	\centering
	\includegraphics[width=0.5\textwidth]{10DES.png}
	\caption{Gráficos de dispersión por pares de indicadores 2020}
	\label{fig:D10}
\end{figure}

El la figura \ref{fig:D10} se observa baja dispersión entre las variables y un considerable grado de traslape.

A continuación se obtuvieron las proyecciones visuales mediante la reducción de características con LDA, KFDA y GNLPDA.

\begin{figure}[htbp]
	\centering
	
	% Primera fila
	\begin{subfigure}[b]{0.3\textwidth}
		\centering
		\includegraphics[width=\textwidth]{mejor_escenario_5.5.indicadores2020_1_mass_2d.png}
	\end{subfigure}
	\hfill
	\begin{subfigure}[b]{0.3\textwidth}
		\centering
		\includegraphics[width=\textwidth]{mejor_escenario_5.5.indicadores2020_2_kfda_2d.png}
	\end{subfigure}
	\hfill
	\begin{subfigure}[b]{0.3\textwidth}
		\centering
		\includegraphics[width=\textwidth]{mejor_escenario_5.5.indicadores2020_3_gnlpda_2d.png}
	\end{subfigure}
	
	\caption{Comparativa geométrica de los espacios de proyección para el conjunto de datos
		Indicadores 2020 (2D).}
	\label{fig:R17}
\end{figure}

\begin{figure}[htbp]
	\centering
	
	% Primera fila
	\begin{subfigure}[b]{0.3\textwidth}
		\centering
		\includegraphics[width=\textwidth]{mejor_escenario_5.5.indicadores2020_1_mass_3d.png}
	\end{subfigure}
	\hfill
	\begin{subfigure}[b]{0.3\textwidth}
		\centering
		\includegraphics[width=\textwidth]{mejor_escenario_5.5.indicadores2020_2_kfda_3d.png}
	\end{subfigure}
	\hfill
	\begin{subfigure}[b]{0.3\textwidth}
		\centering
		\includegraphics[width=\textwidth]{mejor_escenario_5.5.indicadores2020_3_gnlpda_3d.png}
	\end{subfigure}
	
	\caption{Comparativa geométrica de los espacios de proyección para el conjunto de datos
		Indicadores 2020 (3D).}
	\label{fig:R18}
\end{figure}


Las Figuras \ref{fig:R17} y \ref{fig:R18} corroboran el comportamiento geométrico analizado
previamente, extendiendo la evaluación a la actualización socioeconómica del año 2020. La
consistencia geométrica es evidente: el enfoque lineal (LDA MASS) reafirma su capacidad
para estructurar la trayectoria curva que representa la progresión ordinal de los niveles de
rezago, logrando una alta exactitud (0.938). De manera análoga, el mapeo implícito (KFDA)
confirma su deficiencia estructural para este dominio; el kernel gaussiano colapsa
nuevamente la relación secuencial, generando un aglomeramiento caótico que desploma su
capacidad predictiva (0.707). Finalmente, el método propuesto (GNLPDA) valida su
superioridad algorítmica al maximizar el rendimiento general (0.959) a través de la
extracción de 39 términos polinomiales. Las proyecciones, especialmente en su vista tridimensional, demuestran que la búsqueda estocástica no solo preserva, sino que acentúa y clarifica la trayectoria continua entre los cinco niveles categóricos. Esta reproducibilidad geométrica a través de diferentes periodos censales consolida la robustez del GNLPDA para generar espacios de características explícitos y parsimoniosos que respetan la naturaleza
ordinal del problema.

\begin{figure}[htbp]
	\centering
	\includegraphics[width=0.8\textwidth]{Boxplot_5.5.indicadores2020_Accuracy.png}
	\caption{Distribución empírica del rendimiento predictivo para el conjunto de datos Indicadores 2020.}
	\label{fig:R46}
\end{figure}

La Figura \ref{fig:R46} detalla el comportamiento de los clasificadores en términos de exactitud
(accuracy). Los modelos SVM(POL2), SVM(POL3) y SVM(Gauss) alcanzan las medianas
más elevadas, situándose consistentemente por encima de 0.95 con rangos muy compactos. 
El método propuesto (GNLPDA) exhibe una mediana superior a 0.92, mostrando una
distribución concentrada con valores de exactitud competitivos. Por su parte, los
clasificadores LDA y QDA presentan medianas situadas por encima del 0.85 con rangos de
dispersión moderadamente compactos. Finalmente, el clasificador KFDA muestra el
desempeño más bajo, con una mediana de exactitud próxima a 0.68 y una mayor dispersión
en los datos en comparación con el resto de los algoritmos evaluados.

\begin{figure}[H]
	\centering
	\includegraphics[width=0.8\textwidth]{CDplot_Nemenyi_5.5.indicadores2020_long_graficos_todos.png}
	\caption{Análisis de significancia estadística post-hoc mediante distancias críticas de
		Nemenyi para el conjunto de datos Indicadores 2020.}
	\label{fig:R47}
\end{figure}

La prueba de Friedman ($\chi^2= 169.856 ,p<0.001 $ ) ratifica la existencia de
diferencias estadísticamente significativas en el desempeño de los algoritmos.
El diagrama de Nemenyi evidencia una jerarquía donde los modelos SVM(POL2,
rango 1.47), SVM (POL3, rango 2.03) y SVM (Gauss, rango 2.50) integran el clúster de
mayor eficacia estadística, sin diferencias significativas entre ellos. El método propuesto
(GNLPDA, rango 4.15) se posiciona en una jerarquía intermedia, seguido por los modelos
LDA (rango 4.85), QDA (rango 6.00) y KFDA (rango 7.00), los cuales ocupan los rangos
inferiores de la distribución. Las líneas horizontales conectan a los clasificadores que, de
acuerdo con el umbral de Diferencia Crítica (CD), no muestran diferencias estadísticas
significativas.

\begin{figure}[htbp]
	\centering
	\includegraphics[width=0.8\textwidth]{CDplot_Nemenyi_5.5.indicadores2020_long_graficos_campeones.png}
	\caption{Análisis de significancia estadística post-hoc y corrección por enmascaramiento
		mediante distancias críticas de Nemenyi para el conjunto de datos Indicadores 2020.}
	\label{fig:R48}
\end{figure}

Al restringir el análisis exclusivamente a los clasificadores no lineales, el test de
Friedman corrobora la existencia de diferencias estadísticamente significativas
($\chi^2=103.811 ,p<0.001$ ). El diagrama de Nemenyi detalla una jerarquía donde
los modelos SVM(POL2, rango 1.47), SVM(POL3, rango 2.03) y SVM(Gauss, rango 2.50)
conforman el clúster de mayor eficacia estadística, sin diferencias significativas entre ellos.
El método propuesto (GNLPDA, rango 4.00) se posiciona en un nivel intermedio,
presentando una diferencia estadística significativa respecto al grupo líder, pero superando
al clasificador KFDA (rango 5.00), el cual se ubica en el extremo inferior de la jerarquía. Las
líneas horizontales conectan a los grupos de clasificadores que, según el umbral de Diferencia
Crítica (CD), no muestran diferencias estadísticas significativas.

